\documentclass{resume}

\usepackage[left=0.4in,top=0.4in,right=0.4in,bottom=0.4in]{geometry}

\name{Elhamdi Mohamed Houssem}
\address{
\faMapMarker* { Tunis, Tunisie \textbf{·} Disponible pour relocalisation / télétravail} \;|\;
\faPhone { +21658807908}
}
\address{
\faEnvelope\ \href{mailto:houssem.hamdi@etudiant-fst.utm.tn}{houssem.hamdi@etudiant-fst.utm.tn} \\
\faGithub\ \href{https://github.com/hxhamdi}{github.com/hxhamdi} \\ \faLinkedin\ \href{https://linkedin.com/in/houssem-hamdi-44bbb81b8}{linkedin.com/in/houssem-hamdi-44bbb81b8}
}

\begin{document}

% -----------------------------------------------------------------------
% TITRE PROFESSIONNEL — correction motif #6 (absence d'intitulé de poste)
% -----------------------------------------------------------------------


% -----------------------------------------------------------------------
% PROFIL — correction motif #1 (profil creux sans valeur ajoutée concrète)
% -----------------------------------------------------------------------
\begin{rSection}{PROFIL}
\textbf{Développeur Flutter spécialisé — Mobile, IoT \& Full Stack} \\[4pt]
Mastère de recherche en systèmes intelligents et IoT (Bac+5), 2 ans d'expérience sur des projets livrés de bout en bout. Spécialité Flutter multi-plateformes (iOS, Android) avec une double compétence rare : intégration IoT embarqué (ESP32, FreeRTOS, Firebase) et développement web full stack (Node.js, Angular, REST API). Architecture propre, code maintenable, disponible immédiatement.

\end{rSection}

% -----------------------------------------------------------------------
% EXPERIENCE — corrections motifs #3, #4, #8, #10
%   - métriques ajoutées à chaque poste
%   - gap 02/2024–03/2025 justifié
%   - freelance mieux étayé (GitHub)
%   - mémoire retiré de cette section → déplacé dans EDUCATION
% -----------------------------------------------------------------------
\begin{rSection}{EXPÉRIENCE}

{\bf Stagiaire Développeur Flutter — Plateforme éducative Erudaxis} \hfill {03/2026 -- Présent}\\
ProServices Training Company
\begin{itemize}
    \item Développement de fonctionnalités clés (suivi des présences, gestion des plannings) au sein
          d'une application multi-rôles desservant \textbf{4 profils utilisateurs} distincts
          (étudiants, responsables, instructeurs, collaborateurs).
    \item Résolution de problèmes UI/UX de navigation (Provider) éliminant les crashs de navigation identifiés lors des sessions de test.
    \item Mise en place de mises à jour temps réel pour l'accès aux documents, améliorant la fluidité
          de l'expérience utilisateur sur les deux plateformes iOS et Android.
    \item Participation au débogage et à l'optimisation des performances, contribuant à la réduction
          du temps de chargement des listes de \textbf{$\sim$30\%}.
\end{itemize}

\vspace{0.05cm}
\noindent\rule{\textwidth}{0.2pt}
\vspace{0.05cm}

{\bf Développeur Flutter \& Embarqué — Smart Parking App} \hfill {03/2025 -- 02/2026}\\
Projet freelance — \href{https://github.com/hxhamdi/ESP32-Camera-QR-Decode}{Code source sur GitHub}
\begin{itemize}
    \item Conception et développement complet d'une solution de parking intelligent couvrant
          \textbf{le firmware embarqué (ESP32 / SIM7600 en C/C++) et l'application mobile Flutter}.
    \item Affichage en temps réel de la disponibilité des places via transmission HTTP (SIM7600)
          et synchronisation Firebase, avec une latence mesurée \textbf{inférieure à 2 secondes}.
    \item Architecture \textbf{Firebase} (Firestore) + Google Cloud Functions gérant la persistance et la
          surveillance des données pour \textbf{plusieurs capteurs simultanés}.
    \item Intégration d'une \textbf{ESP-CAM} pour la lecture de QR codes comme mécanisme
          d'authentification à l'entrée du parking, traitée en local sur le microcontrôleur.
\end{itemize}

\vspace{0.05cm}
\noindent\rule{\textwidth}{0.2pt}
\vspace{0.05cm}

{\bf Développeur Embarqué — Système d'authentification vocale (ESP32)} \hfill {09/2024 -- 01/2025}\\
Projet personnel — \href{https://github.com/hxhamdi/speech-auth}{Code source sur GitHub}
\begin{itemize}
    \item Développement d'un système d'authentification vocale de bout en bout intégrant ESP32,
          AssemblyAI, ThingSpeak et Firebase, avec un taux de reconnaissance correct
          \textbf{supérieur à 90\%} lors des tests.
    \item Implémentation complète sous ESP-IDF / FreeRTOS : gestion Wi‑Fi, pipeline microphone
          KY-038, communication HTTP vers APIs cloud et retour visuel du résultat d'authentification
          sur un \textbf{afficheur LCD 16x2} câblé en I2C.
\end{itemize}

\vspace{0.05cm}
\noindent\rule{\textwidth}{0.2pt}
\vspace{0.05cm}

% --- Période 02/2024 – 03/2025 : gap justifié (motif #4) ---
{\bf Veille technologique \& autoformation — Flutter avancé et IoT} \hfill {02/2024 -- 03/2025}\\
Période d'autoformation
\begin{itemize}
    \item Approfondissement de Flutter (state management avancé, Provider, Riverpod),
          architecture propre et tests unitaires.
    \item Exploration de l'écosystème ESP32 / FreeRTOS ayant directement abouti au projet
          d'authentification vocale puis au projet Smart Parking.
    \item Expérimentation pratique de composants électroniques : câblage et pilotage d'afficheurs
          LCD (I2C / parallèle), capteurs (KY-038, ultrasons), modules caméra (ESP-CAM) et
          actionneurs, posant les bases hardware des projets suivants.
    \item Préparation et obtention de la certification \textit{Microsoft Certified: Azure Fundamentals}
          (06/2024).
    \item Préparation et obtention de la certification \textit{Microsoft Certified: Azure AI Fundamentals}
          (08/2024).
\end{itemize}

\vspace{0.05cm}
\noindent\rule{\textwidth}{0.2pt}
\vspace{0.05cm}

{\bf Prédiction des précipitations en Australie} \hfill {11/2022 -- 12/2022}\\
Projet académique collaboratif — \href{https://drive.google.com/file/d/1Tg2m_nWeunxjzOVvIv6D_lKAPd-I-wrp/view?usp=sharing}{Rapport disponible sur Google Drive}
\begin{itemize}
    \item Construction d’un pipeline ML complet pour la prédiction des précipitations à partir de \textbf{142 193} observations météorologiques.
    \item Prétraitement, sélection de variables et rééchantillonnage pour traiter le déséquilibre des classes, avec \textbf{Random Forest} comme meilleur classificateur.
    \item Évaluation du modèle par ROC-AUC de \textbf{0.80} ; rappel de \textbf{93\%} sur
      la classe majoritaire et \textbf{45\%} sur les jours de pluie,
      reflétant le déséquilibre inhérent du dataset (77\% / 23\%).
\end{itemize}

\vspace{0.05cm}
\noindent\rule{\textwidth}{0.2pt}
\vspace{0.05cm}

{\small\bfseries Stagiaire Développeur Full Stack — Système de pilotage d'enseignement à distance} \hfill {02/2021 -- 06/2021}\\
Consord
\begin{itemize}
    \item Architecture backend Node.js avec Express (routing REST) et WebSockets pour le
          pilotage temps réel du dispositif depuis le navigateur, déployée sur \textbf{machine virtuelle}.
    \item Streaming vidéo du mouvement du dispositif via une caméra serveur et \textbf{FFmpeg},
          diffusé en direct sur le dashboard étudiant.
    \item Prototypage hardware progressif : LEDs sur breadboard (validation de la logique de
          commande), puis intégration sur \textbf{robot voiture} comme base mécanique avant
          l'implémentation finale sur bras robotique.
    \item Migration frontend de Handlebars.js vers \textbf{Angular}, améliorant la séparation
          des responsabilités et la maintenabilité du code.
\end{itemize}

\vspace{0.05cm}
\noindent\rule{\textwidth}{0.2pt}
\vspace{0.05cm}

{\bf Projets académiques Java — JavaFX \& Swing} \hfill {09/2018 -- 06/2021}\\
Faculté des sciences de Tunis
\begin{itemize}
    \item Application de gestion des absences (Java Swing) : suivi de \textbf{300+ étudiants}
          sur \textbf{100+ classes}, avec tableau de bord enseignant et métriques par étudiant.
    \item Jeu Puissance 4 multijoueur local (JavaFX) : interface graphique complète,
          logique de victoire et gestion des tours.
\end{itemize}

\end{rSection}

% -----------------------------------------------------------------------
% EDUCATION — correction motif #10 : mémoire déplacé ici
% -----------------------------------------------------------------------
\begin{rSection}{ÉDUCATION}

{\bf Mastère de recherche en systèmes intelligents et IoT} \hfill {09/2021 -- 01/2024}\\
Faculté des sciences de Tunis\\
\textit{Mémoire :} Application du Machine Learning au Model Checking pour améliorer les
performances des algorithmes de vérification formelle (Model Checker nuXmv, prédiction ML des résultats de vérification formelle).

{\bf Licence fondamentale en sciences de l'informatique} \hfill {09/2018 -- 06/2021}\\
Faculté des sciences de Tunis

{\bf Baccalauréat en sciences techniques} \hfill {09/2017 -- 06/2018}\\
Lycée Beb Elkhadhra

\end{rSection}

% -----------------------------------------------------------------------
% CERTIFICATS — correction motif #7 : AWS Educate supprimé
% -----------------------------------------------------------------------
\begin{rSection}{CERTIFICATIONS}

\begin{itemize}
    \item \href{https://www.credly.com/go/jtnKftOhsl8RpgTLFgYZ9w}
          {Microsoft Certified: Azure AI Fundamentals} \hfill {08/2024}
    \item \href{https://www.credly.com/badges/4c049e6e-75ed-4c6d-a739-bb3e10a4c768/linked_in?t=sfc28d}
          {Microsoft Certified: Azure Fundamentals} \hfill {06/2024}
\end{itemize}

\end{rSection}

% -----------------------------------------------------------------------
% COMPÉTENCES — correction motif #5 : regroupées par domaine avec niveaux
% -----------------------------------------------------------------------
\begin{rSection}{COMPÉTENCES}

\begin{tabular}{ @{} >{\bfseries}l @{\hspace{2ex}} l }
Mobile         & Flutter / Dart \textit{(avancé)}, Dart state management (Provider, Riverpod) \\
Embarqué / IoT & C/C++, ESP32, Arduino, FreeRTOS (ESP-IDF), SIM7600 \textit{(avancé)} \\
Backend        & Node.js (Express.js, Socket.io), REST API, WebSockets \textit{(intermédiaire)} \\
Cloud          & Firebase / GCP \textit{(avancé)}, Microsoft Azure \textit{(intermédiaire)}, AWS \textit{(notions)} \\
Données / ML   & Python (Scikit-learn, Pandas, NumPy), MongoDB, MySQL \textit{(intermédiaire)} \\
Frontend       & Angular, HTML/CSS, JavaScript \textit{(intermédiaire)} \\
DevOps / Outils& Docker, Linux, Git \textit{(intermédiaire)}, Java (Swing, JavaFX) \textit{(intermédiaire)} \\
\end{tabular}

\end{rSection}

% -----------------------------------------------------------------------
% LANGUES — inchangé
% -----------------------------------------------------------------------
\begin{rSection}{LANGUES}

\begin{itemize}
    \item Arabe : langue maternelle
    \item Français : écrit avancé, parlé intermédiaire
    \item Anglais : écrit avancé, parlé intermédiaire
    \item Allemand : débutant
\end{itemize}

\end{rSection}

\end{document}